\subsection{Local objects}
\label{subsec:laxloc-local}

Let
$P \subset \func(\Delta^{1},\D)$
be a locally full sub-$(\infty,2)$-category. We write the objects of $P$ as
$p\colon x \to y$, and its $1$-morphisms, which are commutative squares, as
$\tau = (a,b)\colon p \to p'$,
drawn with $p$ and $p'$ vertical as in \cref{lem:mates}(1).

\begin{defi}
\label{def:laxloc-local}
An object $z \in \D$ is \emph{$P$-local} if
\begin{enumerate}
\item[(O1)] for every object $p\colon x \to y$ of $P$, the functor
$p^{\natural}\colon \Hom_{\D}(y,z) \to \Hom_{\D}(x,z)$
has a fully faithful left adjoint $(p^{\natural})^{L}$, and
\item[(O2)] for every $1$-morphism
$\tau = (a,b)\colon (p\colon x \to y) \to (p'\colon x' \to y')$
of $P$, the commutative square
\begin{equation*}
\begin{tikzcd}[column sep=large]
\Hom_{\D}(y',z) \ar[r, "p'^{\natural}"] \ar[d, "b^{\natural}"'] &
\Hom_{\D}(x',z) \ar[d, "a^{\natural}"] \ar[dl, Rightarrow, shorten=4mm] \\
\Hom_{\D}(y,z) \ar[r, "p^{\natural}"'] & \Hom_{\D}(x,z)
\end{tikzcd}
\end{equation*}
filled by $\tau$ is left Beck--Chevalley, i.e.\
$(p^{\natural})^{L}a^{\natural} \Rightarrow b^{\natural}(p'^{\natural})^{L}$
is invertible.
\end{enumerate}
A $1$-morphism $g\colon z \to z'$ between $P$-local objects is
\emph{$P$-local} if
\begin{enumerate}
\item[(M)] for every object $p\colon x \to y$ of $P$, the square
\begin{equation*}
\begin{tikzcd}[column sep=large]
\Hom_{\D}(y,z) \ar[r, "p^{\natural}"] \ar[d, "g_{\natural}"'] &
\Hom_{\D}(x,z) \ar[d, "g_{\natural}"] \ar[dl, Rightarrow, shorten=4mm, "="'] \\
\Hom_{\D}(y,z') \ar[r, "p^{\natural}"'] & \Hom_{\D}(x,z')
\end{tikzcd}
\end{equation*}
given by associativity is left Beck--Chevalley, i.e.\
$(p^{\natural})^{L}g_{\natural} \Rightarrow g_{\natural}(p^{\natural})^{L}$
is invertible.
\end{enumerate}
Let
$\Pc_{\emptyset,P} \subset \D$
be the locally full sub-$(\infty,2)$-category of $P$-local objects and
$1$-morphisms.
\end{defi}

By \cref{lem:mate-correspondence}(3), $P$-local $1$-morphisms are closed under
composition and contain the identities, and both classes are closed under
equivalence, so $\Pc_{\emptyset,P}$ exists (\cref{subsec:prelim-conventions}).
We will see in \cref{thm:local-pullback} that it is the case $I = Q =
\emptyset$ of \cref{def:IP-local}, which explains the notation.
If the $1$-morphisms of
$P$ are just the equivalences, the $P$-local objects and $1$-morphisms are the
left Kan injective ones of \cite[Def.~1.2,
Rem.~1.3]{dilibertilobbiasousa2022kz}. We
write $\{\sigma\}$ for the locally full sub-$(\infty,2)$-category of
$\func(\Delta^{1},\D)$ generated by a $1$-morphism or a commutative square
$\sigma$, and $\{s_{j}\}_{j}$, or just $\{s_{j}\}$,
for the one whose objects are the $1$-morphisms
$s_{j}$ and whose $1$-morphisms are the equivalences. Then $\Pc_{\emptyset,P}$
is the intersection of the
$\Pc_{\emptyset,\{p\}}$
and
$\Pc_{\emptyset,\{\tau\}}$
over the objects $p$ and $1$-morphisms $\tau$ of $P$.

As $(p^{\natural})^{L}$ is fully faithful, the counit of
$(p^{\natural})^{L} \dashv p^{\natural}$
is invertible exactly at the objects of its essential image. By
\cref{constr:mates}, the left mates of (O2) and (M) are, at $m$, this counit at
$b^{\natural}(p'^{\natural})^{L}m$,
resp.\ at
$g_{\natural}(p^{\natural})^{L}m$,
composed with invertible $2$-morphisms. So, granting (O1), (O2) says that
$b^{\natural}$ carries the essential image of $(p'^{\natural})^{L}$ into that of
$(p^{\natural})^{L}$, and (M) says that $g_{\natural}$ carries the essential
image of $(p^{\natural})^{L}$ on $\Hom_{\D}(y,z)$ into that on
$\Hom_{\D}(y,z')$.

\begin{lem}
\label{lem:laxloc-closure}
\leavevmode
\begin{enumerate}
\item A left adjoint between $P$-local objects is $P$-local.
\item If $d$ is $P$-local and $i\colon e \to d$ is reflective, then $e$ and
$i^{L}$ are $P$-local.
\item Let $i\colon e \to d$ be reflective and $f\colon e \to e'$ a $1$-morphism
such that $i^{L}$ and $fi^{L}$ are $P$-local. Then $f$ is $P$-local.
\end{enumerate}
\end{lem}

\begin{proof}
(1) Let
$g \dashv g^{R}$
with $g\colon z \to z'$. The square of (M) has the right mate
$p^{\natural}(g^{R})_{\natural} \Rightarrow (g^{R})_{\natural}p^{\natural}$,
which is invertible, being $p^{\natural}$ applied to a triangle identity. So the
square is left Beck--Chevalley by \cref{lem:mates-conjugate}.

(2) Let $p\colon x \to y$ be an object of $P$, $(p^{\natural})^{L}$ the left
adjoint of $p^{\natural}$ on $\Hom_{\D}(-,d)$, and
$\lambda := i^{L}_{\natural}(p^{\natural})^{L}i_{\natural}\colon
\Hom_{\D}(x,e) \to \Hom_{\D}(y,e)$.
As $i_{\natural}$ is fully faithful and commutes with $p^{\natural}$, there
are natural equivalences
\begin{equation*}
\mapp(\lambda m,n) \simeq \mapp\bigl((p^{\natural})^{L}im,in\bigr) \simeq
\mapp(im,inp) \simeq \mapp(m,np) \ ,
\end{equation*}
so
$\lambda \dashv p^{\natural}$
\cite[Prop.~5.2.2.8]{htt}. Under $i_{\natural}$,
its unit at $m$ is the unit of $(p^{\natural})^{L}$ at $im$ followed by the unit
of
$i^{L} \dashv i$
at
$(p^{\natural})^{L}(im) \circ p \simeq im$;
both are invertible, the latter as $im$ lies in the essential image of
$i_{\natural}$. So (O1) holds for $e$. For $\tau = (a,b)$ in $P$, (O2) for $d$
gives
\begin{equation*}
b^{\natural}\lambda' =
i^{L}_{\natural}b^{\natural}(p'^{\natural})^{L}i_{\natural}
\simeq i^{L}_{\natural}(p^{\natural})^{L}a^{\natural}i_{\natural} = \lambda
a^{\natural} \ ,
\end{equation*}
with $\lambda'$ defined like $\lambda$ for $p'$. So $b^{\natural}$ carries the
essential image of $\lambda'$ into that of $\lambda$, which is (O2) for $e$.
Now $i^{L}$ is $P$-local by (1).

(3) The objects $d$, $e$ and $e'$ are $P$-local, being sources and targets of
$P$-local $1$-morphisms. By \cref{lem:mate-correspondence}(3),
the left mate of (M) for $fi^{L}$ is the
pasting of those for $i^{L}$ and $f$. As the former two are invertible, the
left mate for $f$ is invertible after precomposition with $i^{L}_{\natural}$,
which is essentially surjective as
$i^{L}_{\natural}i_{\natural} \simeq \mathrm{id}$.
So $f$ is $P$-local.
\end{proof}

\begin{lem}
\label{lem:laxloc-repr}
A $1$-morphism $s$ of $\D$ is coreflective if and only if
$\Pc_{\emptyset,\{s\}} = \D$.
A commutative square $\sigma$ between coreflective $1$-morphisms, drawn with
these vertical, is right Beck--Chevalley if and only if
$\Pc_{\emptyset,\{\sigma\}} = \D$.
\end{lem}

\begin{proof}
The enriched Yoneda embedding
$h\colon \D\opcat \to \func(\D,\Cat)$,
$d \mapsto \Hom_{\D}(d,-)$,
is fully faithful \cite{hinich2020yoneda,heine2023enriched}. It takes an
adjunction $l \dashv r$ of $\D$ to
$r^{\natural} \dashv l^{\natural}$
with the same unit and counit, and every adjunction between representables
arises this way, as adjunctions are detected in the homotopy $2$-category
(\cref{subsec:prelim-adjunctions}). So $s\colon x \to y$ is coreflective if and
only if
$s^{\natural}\colon h(y) \to h(x)$
is a right adjoint with invertible unit. By
\cite[Cor.~5.2.12]{abellangagnahaugseng2024straightening}, $s^{\natural}$ is a
right adjoint if and only if its components are and its transposed naturality
squares, the squares of (M), are left Beck--Chevalley; and a $2$-morphism of
$\func(\D,\Cat)$ is invertible if and only if its components are. This proves
the first claim. For
$\sigma = (a,b)\colon s \to s'$,
$h$ takes the right mate
$as^{R} \Rightarrow s'^{R}b$
to the left mate
$(s^{\natural})^{L}a^{\natural} \Rightarrow b^{\natural}(s'^{\natural})^{L}$
of $h(\sigma)$, whose components are the $2$-morphisms of (O2), and $h$
reflects invertibility.
\end{proof}

\begin{prop}
\label{prop:laxloc-adjunction}
Let $s$ be a $1$-morphism of $\D$. Then $Ls$ is coreflective if and only if
every object $Rc$ and every $1$-morphism $Rg$ is $\{s\}$-local. If $\sigma$ is
a commutative square between such $1$-morphisms, then $L\sigma$ is right
Beck--Chevalley if and only if every $Rc$ is $\{\sigma\}$-local.
\end{prop}

\begin{proof}
By naturality of \eqref{eq:laxloc-adjunction}, the squares of
\cref{def:laxloc-local} for $Ls$, $L\sigma$ and an object $c$ or a $1$-morphism
$g$ of $\C$ are equivalent to those for $s$, $\sigma$ and $Rc$, $Rg$. So this
is \cref{lem:laxloc-repr} in $\C$.
\end{proof}
